 %%%%%%%%%%%%%%%%%%%%%%%%%%%%%%%%%%%%%%%%%
% Beamer Presentation
% LaTeX Template
% Version 1.0 (10/11/12)
%
% This template has been downloaded from:
% http://www.LaTeXTemplates.com
%
% License:
% CC BY-NC-SA 3.0 (http://creativecommons.org/licenses/by-nc-sa/3.0/)
%
%%%%%%%%%%%%%%%%%%%%%%%%%%%%%%%%%%%%%%%%%

%----------------------------------------------------------------------------------------
%	PACKAGES AND THEMES
%----------------------------------------------------------------------------------------

\documentclass[9pt,aspectratio=169]{beamer}
\usepackage[utf8]{inputenc}
\usepackage{bm}
\usepackage{bbm}
\usepackage{graphicx}
\usepackage{cmbright}
\usefonttheme{professionalfonts}
%\usefonttheme{serif}
% The Beamer class comes with a number of default slide themes
% which change the colors and layouts of slides. Below this is a list
% of all the themes, uncomment each in turn to see what they look like.

%\usetheme{default}
%\usetheme{AnnArbor}
% \usetheme{Antibes}
%\usetheme{Bergen}
%\usetheme{Berkeley}
%\usetheme{Berlin}
%\usetheme{Boadilla}
% \usetheme{CambridgeUS}
%\usetheme{Copenhagen}
%\usetheme{Darmstadt}
% \usetheme{Dresden}
\usetheme{Frankfurt}
%\usetheme{Goettingen}
% \usetheme{Hannover}
% \usetheme{Ilmenau}
% \usetheme{JuanLesPins}
%\usetheme{Luebeck}
%\usetheme{Madrid}
%\usetheme{Malmoe}
%\usetheme{Marburg}
% \usetheme{Montpellier}
% \usetheme[width=2.5cm]{PaloAlto}
%\usetheme{Pittsburgh}
% \usetheme{Rochester}
%\usetheme{Singapore}
% \usetheme{Szeged}
% \usetheme{Warsaw}
% \setbeamertemplate{title page}[default][colsep=-4bp,rounded=false]
\setbeamertemplate{headline}{}
\defbeamertemplate{footline}{myframe number}
{
  \hspace{1pt}
    \usebeamercolor[structure]{page number in head/foot}%
  \usebeamerfont{page number in head/foot}%
  \raisebox{1.7ex}[0pt][0pt]{% <--- change here
    \insertframenumber\,/\,\inserttotalframenumber\kern1em}%
}
\setbeamertemplate{footline}[myframe number]
\AtBeginSection[]
{
  \begin{frame}
    \frametitle{Outline}
    \begin{minipage}{1.0\linewidth}
      \tableofcontents[currentsection,currentsubsection]
    \end{minipage}
  \end{frame}
}

\AtBeginSubsection[]
{
  \begin{frame}
    \frametitle{Outline}
    \begin{minipage}{1.0\linewidth}
      \tableofcontents[currentsection,currentsubsection]
    \end{minipage}
  \end{frame}
}

% As well as themes, the Beamer class has a number of color themes
% for any slide theme. Uncomment each of these in turn to see how it
% changes the colors of your current slide theme.
%\definecolor{mycolor}{rgb}{.8,.0,.3}
%\setbeamercolor*{palette primary}{use=structure,fg=white,bg=mycolor}
%\usecolortheme{albatross}
%\usecolortheme{beaver}
%\usecolortheme{beetle}
%\usecolortheme{crane}
%\usecolortheme{dolphin}
%\usecolortheme{dove}
%\usecolortheme{fly}
%\usecolortheme{lily}
%\usecolortheme{orchid}
% \usecolortheme{rose}
%\usecolortheme{seagull}
\usecolortheme{seahorse}
%\usecolortheme{whale}
%\usecolortheme{wolverine}
%\usecolortheme[named=mycolor]{structure}
%\setbeamertemplate{footline} % To remove the footer line in all slides uncomment this line
%\setbeamertemplate{footline}[page number] % To replace the footer line in all slides with a simple slide count uncomment this line

%\setbeamertemplate{navigation symbols}{} % To remove the navigation symbols from the bottom of all slides uncomment this line


\usepackage{booktabs} % Allows the use of \toprule, \midrule and \bottomrule in tables

%----------------------------------------------------------------------------------------
%	TITLE PAGE
%----------------------------------------------------------------------------------------

\title[Lecture 5: Constrained Optimization II]{Lecture 5: Constrained
  Optimization II}

\author{Junnan Zhang} % Your name
\institute[Paula and Gregory Chow Institute for Studies in Economics\\
Xiamen University] % Your institution as it will appear on the bottom of every slide, may be shorthand to save space
{
	Paula and Gregory Chow Institute for Studies in Economics \\ Xiamen University  \\ % Your institution for the title page
	\medskip
    \textit{Slides Prepared by Xiaoling Mei}
	%\textit{luciamay@smith.com} % Your email address
}
\date{Fall, 2026} % Date, can be changed to a custom date
\setbeamertemplate{itemize items}[default]
\setbeamertemplate{sidebar}[sidebar right]
% \definecolor{myco}{HTML}{a67678}
\definecolor{myco}{HTML}{8776a6}
\setbeamercolor{itemize item}{fg=myco} % all frames will have red bullets
\setbeamercolor{itemize subitem}{fg=myco} % all frames will have red bullets
\setbeamercolor{block title}{bg=myco}
\setbeamercolor{structure}{fg=myco}
% \setbeamercovered{transparent}
\newcommand{\RR}{\mathbbm R}

\begin{document}



\begin{frame}
\titlepage % Print the title page as the first slide
\end{frame}


\begin{frame}
    \frametitle{Introduction}
    This lecture focuses on two other aspects of the Lagrangian approach:
    \begin{itemize}
        \item the sensitivity of the optimal value of the objective function
        to changes in the parameters
        \item the constraint qualifications that are a subtle but necessary
        hypothesis in the Lagrangian approach
    \end{itemize}
\end{frame}		

\section{The Meaning of the Multiplier}

\begin{frame}
	\frametitle{The Meaning of the Multiplier}
    Consider the simplest problem with one equality constraint:
    \begin{align}
        \text{max}\ \ & f(x,y) \label{eq1} \\
        \text{s.t.}\ \ & h(x,y) = a \label{eq2} 
    \end{align}
\end{frame}	
				
\begin{frame}
	\frametitle{The Meaning of the Multiplier}
	\begin{block}{\textbf{Theorem 19.1}}
		\begin{itemize}
            \item Let $f$ and $h$ be $C^1$ functions of two variables. For
            any fixed value of the parameter $a$, let $(x^*(a),y^*(a))$ be
            the solution of problem~\eqref{eq1}--\eqref{eq2} with
            corresponding multiplier $\mu^*(a)$.
            \item Suppose that $x^*,y^*$ and $\mu^*$ are $C^1$ functions of $a$
            and that NDCQ holds at $(x^*(a),y^*(a))$. Then
            $$\mu^*(a) = \frac{d}{da}f(x^*(a),y^*(a))$$
        \end{itemize} 
	\end{block}
\end{frame}
		
						
\begin{frame}
	\frametitle{Example 19.1} \textbf{Example 19.1}: Previously, we found
    that a maximizer of $f(x_1,x_2) = x_1^2x_2$ on the constraint set
    $2x_1^2+x_2^2 = 3$ is $x_1 = x_2 = 1$, with multiplier $\mu = 0.5$.
	\begin{itemize}
		\item The maximum value of $f$ is $f^* = f(1,1) = 1$
		\item Redo the problem, this time using the constraint
        $$2x_1^2+x_2^2 = 3.3$$
		\item The same computation yields the solution $x_1 = x_2 =
        \sqrt{1.1}$, with maximum value $f^* = (1.1)^{3/2} \approx 1.1537$,
        an increase of 0.1537 over the original $f^*$
	\end{itemize}
\end{frame}			
		
		
\begin{frame}
	\frametitle{Example 19.1}
	\textbf{Example 19.1}: 
	\begin{itemize}
		\item On the other hand, Theorem 19.1 predicts that changing the
        right-hand side of the constraint by 0.3 unit would change the
        maximum value of the objective function by roughly
		$$0.3\mu = 0.3*0.5 = 0.15 \ \text{unit},$$
		an approximation correct to two decimal places
	\end{itemize}
\end{frame}		
		


\begin{frame}
	\frametitle{Several Equality Constraints: Theorem}
	\begin{block}{\textbf{Theorem 19.2}}
        Let $f,h_1,\cdots,h_m$ be $C^1$ functions on $\RR^n$. Let $\mathbf{a}
        = (a_1,\cdots,a_m)$ be an $m$-tuple of exogenous parameters, and
        consider the problem of maximizing $f(x_1,\cdots,x_n)$ subject to
        \begin{equation*}
            h_1(x_1,\cdots,x_n)= a_1, \ldots, h_m(x_1,\cdots,x_n) = a_m.
        \end{equation*}
        Let $x_1^*(\mathbf{a}), \cdots, x_n^*(\mathbf{a})$ denote the
        solution of the problem with corresponding Lagrange multipliers
        $\mu_1^*(\mathbf{a}),\cdots, \mu_m^*(\mathbf{a})$. Suppose further
        that the $x_i^*$'s and $\mu_j^*$'s are differentiable functions of
        $(a_1,\cdots,a_m)$ and that NDCQ holds. Then, for each $j = 1,\cdots,
        m,$
        \begin{equation*}
            \mu_j^*(a_1,\cdots,a_m) = \frac{\partial}{\partial
              a_j}f(x_1^*(a_1,\cdots,a_m),\cdots,x_n^*(a_1,\cdots,a_m))
        \end{equation*}
	\end{block}
\end{frame}

\begin{frame}
	\frametitle{Several Inequality Constraints}
	\begin{block}{\textbf{Theorem 19.3}}
        Let $\mathbf{a^*} = (a_1^*,\cdots,a_k^*)$ be a $k$-tuple. Consider
        the problem of maximizing $f(x_1, \cdots, x_n)$ subject to
        \begin{equation*}
            g_1(x_1,\cdots,x_n) \leq a_1^*, \ldots, g_k(x_1,\cdots,x_n) \leq
            a_k^*.
        \end{equation*}
        Let $x_1^*(\mathbf{a^*}), \cdots, x_n^*(\mathbf{a^*})$ denote the
        solution of the problem with corresponding Lagrange multipliers
        $\lambda_1^*(\mathbf{a^*}),\cdots, \lambda_k^*(\mathbf{a^*}).$
        Suppose that as $\mathbf{a}$ varies near $\mathbf{a^*}$,
        $x_1^*,\cdots,x_n^*$ and $\lambda_1^*,\cdots,\lambda_k^*$ are
        differentiable functions of $(a_1,\cdots,a_k)$ and that NDCQ holds at
        $\mathbf{a}^*$. Then, for each $j = 1,\cdots, k,$
        \begin{equation*}
            \lambda_j^*(a_1^*,\cdots,a_k^*) = \frac{\partial}{\partial
              a_j}f\left(x_1^*(a_1^*,\cdots,a_k^*),\cdots,x_n^*(a_1^*,
                \cdots, a_k^*)\right)
        \end{equation*}
	\end{block}
\end{frame}	

\begin{frame}
	\frametitle{Interpreting the Multiplier} 
    
	Think of the objective function as the profit function of a firm and
    think of $a_j$ as the amount of available input in the production
    process. Then $\lambda_j = \partial f/\partial a_j$ is called the
    \emph{internal value}, or \emph{shadow price} of input $j$:
	\begin{itemize}
		\item It tells how valuable another unit of certain input would be
        to the firm's profits
		\item It tells the maximum amount the firm would be willing to pay to
        acquire another unit of the input
	\end{itemize}

\end{frame}	




\begin{frame}
    \frametitle{Examples}

    \textbf{Example 19.2} Previously, we computed that the maximizer of $xyz$
    on the constraint set
    $$x+y+z\leq 1,\, x\geq 0,\, y\geq 0,\, z\geq 0$$
    is $x = y =z =1/3$, where $xyz = 1/27$. The four multipliers are
    $1/9, 0, 0$ and $0$ respectively.
    \begin{itemize}
        \item[(a)] If we change the first constraint to $x+y+z\leq0.9$, we
        compute that the solution occurs at $x = y = z = 0.3$, where $xyz =
        0.027$. Theorem 19.3 predicts that the new optimal value would be
        $\frac{1}{27} + \frac{1}{9} \cdot \left(-\frac{1}{10}\right) \approx
        0.0259$, an estimate that is off by only 0.0011 or 4\%.
        \item[(b)] If, instead, we change the second constraint form $x\geq0$
        to $x\geq0.1$, we do not change the solution or the optimum value
        because the new region is a subset of the old region and it still
        contains the optimal point for the old region. This result is
        consistent with Theorem 19.3 since the multiplier for the
        (nonbinding) constraint $x\geq0$ was zero.
    \end{itemize}
\end{frame}	

\section{Envelope Theorems}
\begin{frame}
	\frametitle{Envelope Theorems}
	\begin{itemize}
		\item Previous theorems are special cases of a class of theorems
        which describe how the optimal value of the objective function in a
        parameterized optimization problem changes as one of the parameter
        changes.
		\item Such theorems are called \textbf{Envelope theorems}.
	\end{itemize}
\end{frame}	


\begin{frame}
	\frametitle{Envelope Theorems}
	\begin{block}{\textbf{Theorem 19.4}}
        Let $f(\mathbf{x}; a)$ be a $C^1$ function of $\mathbf{x}\in \RR^n$
        and the scalar $a$. For each choice of the parameter $a$, consider
        the unconstrained maximization problem: $\max_{\mathbf x}
        f(\mathbf{x}; a)$. Let $\mathbf{x^*}(a)$ be a solution of this
        problem. Suppose that $\mathbf{x^*}(a)$ is a $C^1$ function of $a$.
        Then
        \begin{equation*}
            \frac{d}{da}f(\mathbf{x^*}(a); a) = \frac{\partial}{\partial
              a}f(\mathbf{x^*(a); a})
        \end{equation*}
	\end{block}
\end{frame}	


\begin{frame}
	\frametitle{Envelope Theorems: Examples}
	\textbf{Example 19.3} Consider the problem 
	$$\max_x\ \ f(x,a) = -a^3x^4+15x^3-e^ax^2+17$$
	around $a = 1$. Since $f$ is a quartic polynomial in $x$ with a negative
    leading coefficient, when $a = 1$, $f\rightarrow -\infty$ as
    $x\rightarrow \pm \infty$. So $f$ admits a a finite global maximizer
    $x^*(a)$ for each value of $a$ near 1. Hence,
    $$\frac{d}{da}f(x^*(a),a) = \frac{\partial}{\partial a}f(x^*(a),a) =
    -3a^2x^{*4}-e^ax^{*2}<0$$ Without solving for the optimal $x^*(a)$, we
    can tell that as $a$ increases beyond 1, $f(x^*(a),a)$ is a decreasing
    function of $a$. The peak of the graph of the function decreases as $a$
    increases.
	
\end{frame}


\begin{frame}
	\frametitle{Envelope Theorems: Examples} \textbf{Example 19.4} What will
    be the effect of a unit increase in $a$ on the value of
    $$\max_x\ \ f(x,a) = -x^2+2ax+4a^2$$ for each $a$? Since $f'(x) = -2x+2a
    = 0$, $x^*(a) = a$. Then
    $$f(x^*(a),a) = f(a,a) = -a^2+2a*a+4a^2 = 5a^2,$$
    which will increase at a rate of $10a$ as $a$ increases. Instead, if we
    apply the envelope theorem,
    $$\frac{df^*}{da} = \frac{\partial f}{\partial a}(x^*(a),a) = 2x+8a =
    10a$$ since $x^*(a) = a$.
\end{frame}

\begin{frame}
	\frametitle{Envelope Theorems: Examples}
	\textbf{Example 19.5}
	\begin{itemize}
		\item A silicon Valley firm produces an output of microchips denoted
        by $y$ and has a cost function $c(y)$ with $c'(y)>0$ and $c''(y)>0$.
		\item Of the chips it produces, a fraction $1-\alpha$ are unavoidably defective and cannot be sold. 
		\item Working chips can be sold at price $p$, and
		the microchip market is highly competitive.
	\end{itemize}
	How will an increase in production quality affect the firm’s profit?
\end{frame}


\begin{frame}
	\frametitle{Envelope Theorems: Examples} \textbf{Example 19.5} The firm’s
    profit function is given by
    $$\pi(p,\alpha) = \max_y\ \ [p\alpha y-c(y)]$$\vspace{-1em}
	\begin{itemize}
		\item The conditions on the cost function guarantee that there is a
        nonzero profit-maximizing output $y^*(\alpha)$ which depends smoothly
        on $\alpha$.
		\item The derivative of optimal profit $\pi$ with respect to $\alpha$ is 
		$$\frac{d\pi}{d\alpha} = \frac{\partial}{\partial \alpha}(p\alpha y-c(y))
        = py >0$$ 
	\end{itemize}
\end{frame}


\begin{frame}
	\frametitle{Constrained Problems: Theorem}
    \begin{block}{\textbf{Theorem 19.5}}
        Let $f,h_1,\cdots,h_k$: $\RR^n \times \RR^1\rightarrow \RR^1$ be $C^1$
        functions. Let $\mathbf{x^*}(a) = (x_1^*(a),\cdots,x_n^*(a))$ denote
        the solution of the problem of maximizing $\mathbf{x}\rightarrow
        f(\mathbf{x},a)$ on the constraint set
        $$h_1(\mathbf{x}; a) = 0,\cdots, h_k(\mathbf{x}; a) = 0$$
        for any fixed choice of the parameter $a$. Suppose that
        $\mathbf{x^*}(a)$ and the Lagrangian multipliers
        $\mu_1(a),\cdots,\mu_k(a)$ are $C^1$ functions of $a$ and that the
        NDCQ holds. Then,
        $$\frac{d}{da}f(\mathbf{x^*}(a); a) = \frac{\partial L}{\partial
          a}(\mathbf{x^*}(a),\mu(a); a).$$
    \end{block}
	
\end{frame}	


\begin{frame}
	\frametitle{Envelope Theorems: Example} \textbf{Example 19.6} Consider
    the problem
	\begin{align}
	\max \quad f(x,y) & = xy \nonumber \\ 
	\text{s.t.}\quad x^2+y^2 &\leq 1 \nonumber 
	\end{align}
    Now change the constraint to $x^2+1.1y^2 \leq 1$. If we write the
    constraint as $x^2+ay^2 \leq 1$, the Lagrangian for the parameterized
    problem is $L(x,y,\lambda;a) = xy-\lambda(x^2+ay^2-1)$.
\end{frame}


\begin{frame}
    \frametitle{Envelope Theorems: Example}
    \begin{itemize}
        \item The solution for the original $(a=1)$ problem was
        $x=y=\frac{1}{\sqrt{2}}$, $\lambda = \frac{1}{2}$
        \item The Envelope Theorem tells us that as a changes from 1 to 1.1,
        the optimal value of f changes by approximately $\frac{\partial
          L}{\partial
          a}\left(\frac{1}{\sqrt{2}},\frac{1}{\sqrt{2}},\frac{1}{2};1\right)*0.1$
        \item Since $\frac{\partial L}{\partial a} = -\lambda y^2 = -1/4$,
        the optimal value will decrease by approximately $0.1*(1/4) = 0.025$
        to 0.475
        \item Calculate directly that the solution to the new problem is
        $x=\frac{1}{\sqrt{2}}, y=\frac{1}{\sqrt{2.2}}$, with $f^*\approx
        0.4767$.
    \end{itemize}

\end{frame}

\section{Smooth Dependence on the Parameters}

\begin{frame}
    \frametitle{Smooth Dependence on the Parameters}
    \begin{block}{\bfseries Theorem 19.9}
        Consider the problem of maximizing $f(x; a)$ subject to $h_1(x; a)=0$,
        $h_2(x; a)=0$, \ldots, $h_k(x; a)=0$. Let $x^*(a)$ be the solution of the
        parameterized constrained maximization problem and let $\mu^*(a)$ be
        the corresponding Lagrange multiplier. If the Hessian matrix of the
        Lagrangian is nonsingular at the point $(x^*(a_0), \mu^*(a_0); a_0)$,
        then:
        \begin{enumerate}
            \item $x^*(a)$ and $\mu^*(a)$ are $C^1$ functions of $a$ at $a =
            a_0$; and
            \item the NDCQ holds at $(x^*(a_0), \mu^*(a_0); a_0)$.
        \end{enumerate}
    \end{block}
\end{frame}

\section{Constraint Qualifications}

\begin{frame}{Constraint Qualifications}
    \begin{itemize}
        \item In Lecture 4, we talked about NDCQ.
        \item There are theorems that do not impose constraint qualifications.
    \end{itemize}
\end{frame}

\begin{frame}
    \frametitle{Constraint Qualifications: Theorem}
    \begin{block}{\textbf{Theorem 19.10}}
        Let $f$ and $h$ be $C^1$ functions of two variables. Suppose that
        $\mathbf{x}^*=(x_1^*,x_2^*)$ is a solution of the problem of
        maximizing $f(x_1,x_2)$ subject to $\{(x_1,x_2):h(x_1,x_2)=c\}$.

        Construct the Lagrangian $L(x_1,x_2,\mu_0,\mu_1)
        =\mu_0f(x_1,x_2)-\mu_1[h(x_1,x_2)-c]$, including a multiplier $\mu_0$
        for the objective function. Then, there exist multipliers $\mu_0^*$
        and $\mu_1^*$ such that:
        \begin{enumerate}
            \item[(a)] $\mu_0^*$ and $\mu_1^*$ are not both zero,
            \item[(b)] $\mu_0^*$ is either 0 or 1, and
            \item[(c)] the quadruple $(x_1^*,x_2^*,\mu_0^*,\mu_1^*)$
            satisfies the equations
        \end{enumerate}
        \begin{align*}
            \frac{\partial L}{\partial x_1}
            &=\mu_0^*\frac{\partial f}{\partial x_1}(x_1^*,x_2^*)
              -\mu_1^*\frac{\partial h}{\partial x_1}(x_1^*,x_2^*)=0,\\
            \frac{\partial L}{\partial x_2}
            &=\mu_0^*\frac{\partial f}{\partial x_2}(x_1^*,x_2^*)
              -\mu_1^*\frac{\partial h}{\partial x_2}(x_1^*,x_2^*)=0,\\
            \frac{\partial L}{\partial \mu_1}
            &=c-h(x_1^*,x_2^*)=0.
        \end{align*}

    \end{block}
\end{frame}

\begin{frame}
    \frametitle{Constraint Qualifications: Theorem}
    \begin{block}{\textbf{Theorem 19.11}}
        Suppose that $f,g_1,\ldots,g_k$ are $C^1$ functions of $n$ variables.
        Suppose that $\mathbf{x}^*$ is a local maximizer of $f$ on the
        constraint set defined by the $k$ inequalities
        $g_1(x_1,\ldots,x_n)\leq b_1,\ldots, g_k(x_1,\ldots,x_n)\leq b_k$.

        Form the Lagrangian
        \begin{equation*}
            L(x_1,\ldots,x_n,\lambda_0,\lambda_1,\ldots,\lambda_k)
            ={}\lambda_0f(\mathbf{x})-\lambda_1[g_1(\mathbf{x})-b_1]
            -\cdots -\lambda_k[g_k(\mathbf{x})-b_k],
        \end{equation*}
        with a multiplier $\lambda_0$ for the objective function. Then,
        there exist multipliers
        $\lambda^*=(\lambda_0^*,\lambda_1^*,\ldots,\lambda_k^*)$ such that:
        \begin{enumerate}
            \item[(a)] $\dfrac{\partial L}{\partial x_1}
            (\mathbf{x}^*,\lambda^*)=0,\ldots,
            \dfrac{\partial L}{\partial x_n}
            (\mathbf{x}^*,\lambda^*)=0$,
            \item[(b)] $\lambda_1^*[g_1(\mathbf{x}^*)-b_1]=0,\ldots,
            \lambda_k^*[g_k(\mathbf{x}^*)-b_k]=0$,
            \item[(c)] $\lambda_1^*\geq0,\ldots,\lambda_k^*\geq0$,
            \item[(d)] $g_1(\mathbf{x}^*)\leq b_1,\ldots,
            g_k(\mathbf{x}^*)\leq b_k$,
            \item[(e)] $\lambda_0^*=0$ or $1$, and
            \item[(f)] $(\lambda_0^*,\lambda_1^*,\ldots,\lambda_k^*)
            \neq(0,0,\ldots,0)$.
        \end{enumerate}
    \end{block}
\end{frame}

%
%\begin{frame}
%\frametitle{Homework}
%\begin{itemize}
%	\item Exercise 19.3
%	\item Exercise 19.12
%	\item Exercise 19.14(only check that in Exercise 18.7)
%\end{itemize}
%\end{frame}


\end{document}
